% ============================================================
%  1. fejezet, Bevezetés
% ============================================================
\chapter{Bevezetés}
\label{ch:intro}

Az idősor-előrejelzés a gépi tanulás egyik kiterjedten kutatott
részterülete, amelynek gyakorlati jelentőségét a valós adatok szinte
kivétel nélkül jelen lévő zajossága adja. Szenzorkiesések, aggregációs
torzítások és piaci mikrostruktúra-zaj egyaránt rontják a jel/zaj arányt,
amelyre a kurrens modellek érzékenynek bizonyulnak.
A hosszú távú idősor-előrejelzés (LTSF) feladata ennél is nehezebb:
adott $L$ hosszúságú előzmény alapján kell becsülni a következő $H$ lépést,
ahol $H$ nagyságrendileg elérheti vagy meghaladhatja $L$-et.

\section{Motiváció}

A modern szekvencia-model-családok, a Mamba-alapú szelektív állapottér-modellek
~\cite{gu2021s4,gu2023mamba} és a klasszikus rekurrens hálózatok~\cite{hochreiter1997lstm}
az adatminőségi problémát implicit általánosítás révén kezelik: a zaj- és
jelkomponens szétválasztását kizárólag az adatvezérelt reprezentációtanulásra
bízzák.
Zajmentes adatokon ez az eljárás kielégítően teljesít, ám a megnövelt
zajszintre való általánosítóképesség empirikusan nem garantált.

A Kalman-szűrő~\cite{kalman1960new} ezzel szemben \emph{explicit}
statisztikai modellt állít fel a folyamatzajra ($Q$) és a mérési zajra ($R$),
és ezek arányából származtat egy $K$ gain együtthatót, amely az
előzmény-állapot és az aktuális megfigyelés közötti optimális súlyozást
határozza meg. Ennek az induktív torzításnak egy differenciálható
neurális hálózatba történő beépítése elméleti szempontból ígéretes
megközelítés a zajrobusztusság növelésére.

Ezen megfontolásra épül a jelen szakdolgozatban bemutatott \textbf{KCMamba}
architektúra.

\section{Célkitűzések}

Három konkrét kérdést vizsgálunk:

\begin{enumerate}
    \item Versenyképes-e a KCMamba tiszta adatokon a standard LTSF
          modellekkel szemben?
    \item Mekkora MSE-degradációt szenvednek el a modellek szintetikusan
          zajosított adatokon ($\sigma \in \{0, 1, 3, 5\}$)?
    \item Kevés tréning-ablak (stride $\in \{1, 16\}$) esetén melyik
          architektúra generalizál jobban?
\end{enumerate}

\section{Hozzájárulások}

\begin{itemize}
    \item A \textbf{KCMamba} architektúra leírása: tanult $Q_\text{net}$,
          $R_\text{net}$, $K_\text{net}$ alhálózatok differenciálható
          Kalman-szűrőn belül, párhuzamos prefix-scan GPU-implementációval.
    \item Mérések arról, hogy a tanult $K$ valóban adaptívan változik a
          zajszinttel: $\sigma=0$-nál $K \approx 0{,}71$ (passzív átengedés),
          $\sigma=5$-nél $K \approx 0{,}07$ (erős szűrés).
    \item Összehasonlítás három adatkészleten (Exchange-Rate, ETTh1, ETTh2):
          a Fair Mamba átlag $+28\,234\%$ MSE-degradációt mutat zajjal,
          a KCMamba $+1\,528\%$-ot.
    \item Nyílt forráskódú implementáció (Python/PyTorch/CUDA) és egy
          aszinkron kereskedési rendszer a modell deployálásához.
\end{itemize}

\section{Dolgozat felépítése}

A~\ref{ch:irodalom}. fejezet összefoglalja a Kalman-szűrő, az SSM-ek és
az LTSF-kutatás releváns irodalmát.
A~\ref{ch:architektura}. fejezet ismerteti a KCMamba architektúrát
és matematikai hátterét.
A~\ref{ch:kiserletek}. fejezetben a kísérleti beállítás és az eredmények
szerepelnek.
A~\ref{ch:implementacio}. fejezet a megvalósított szoftverrendszert
tárgyalja.
Az~\ref{ch:osszegzes}. fejezet levonja a következtetéseket.
